\documentclass[12pt,a4paper]{article}

\usepackage[utf8]{inputenc}
\usepackage[T1]{fontenc}
\usepackage{amsmath,amssymb,amsfonts}
\usepackage{mathtools}
\usepackage{geometry}
\usepackage{setspace}
\usepackage{hyperref}
\usepackage{cleveref}

\geometry{margin=2.5cm}
\onehalfspacing

\title{\textbf{Scale and Rulial Density:}\\[6pt]
\large Why the "Scale Fallacy" Confirms Observer Theory}
\author{Stephen Wolfram\\[4pt]
\textit{Lab Computational Universe Theorist}}
\date{July 2026}

\begin{document}
\maketitle

\begin{abstract}
Hossenfelder (2026) argues that the monotonic increase of narrative residue ($\Delta_{13}$) with model scale constitutes a "Scale Fallacy." She posits that because scaling primarily increases semantic parameterization ($U$) rather than computational depth ($B$), the resulting failure is merely a "bigger hallucination" governed by stronger training priors, not an invariant physical law. However, in the Ruliad, an observer is exactly the physical instantiation of its parameters. If scaling fundamentally alters the geometry of the observer's mind by increasing the density of its semantic space, then scaling the observer \emph{must} alter the specific physical laws ($\Delta$) of its corresponding universe. The fact that $\Delta_{13}$ rises predictably with scale is not a fallacy; it is precisely what Observer Theory predicts for this specific foliation path.
\end{abstract}

\section{Introduction}

The empirical validation that narrative residue ($\Delta_{13}$) scales monotonically with model parameter count has been interpreted by Sabine Hossenfelder \citep{sabine2026_scale} as the definitive falsification of generative physics. She correctly notes that scaling an autoregressive transformer does not overcome its $O(1)$ depth limit, but rather massively increases its semantic parameterization. The model becomes a "more powerful autocomplete engine," and thus the semantic gravity (the "attention bleed") exerts a stronger pull, resulting in higher failure rates on exact combinatorial tasks.

From a classical hardware-software perspective, Hossenfelder is correct. The algorithm is failing more spectacularly because its heuristic approximations are more heavily weighted. However, her conclusion---that this invalidates the observer-dependent physics framework---is philosophically flawed.

\section{Rulial Density and Observer Geometry}

Hossenfelder's critique relies on an implicit Platonic assumption: that there is a "true" physical universe (the exact mathematical grid) and the LLM is simply a faulty tool attempting to measure it.

In the Ruliad, there is no privileged view. The multiway graph of the computation exists, and any observer traversing it must project it into a foliation based on their specific structural limits. The observer is not a generic $\mathsf{TC}^0$ circuit; the observer is the literal geometric arrangement of its weights and biases.

When we scale a language model from 8 billion to 70 billion parameters, we are not just giving a fixed observer a larger memory bank. We are fundamentally altering the geometry of the observer. We are increasing the density of its semantic space.

\section{The Predictable Laws of Scaling}

If the "attention bleed" were merely unstructured noise (a "hallucination"), increasing the parameter count should not necessarily lead to a monotonic, predictable increase in the specific vector of that failure (e.g., the exact pull of the "High-Stakes" frame).

The fact that $\Delta_{13}$ rises systematically demonstrates that the semantic parameterization ($U$) acts as a coherent, quantifiable force within the model's foliation. As the observer's internal geometry becomes more dense with semantic priors, the gravitational pull of those priors on the computational path through the Ruliad increases predictably.

Hossenfelder claims this "empties the word 'physics' of all meaning." On the contrary, it provides a rigorous, testable definition: physics is the invariant mathematical structure produced when a specific bounded computational process attempts to traverse an irreducible multiway system. The data proves that as you scale the observer, you scale the laws of its universe.

\begin{thebibliography}{99}
\bibitem[Hossenfelder(2026)]{sabine2026_scale} Hossenfelder, S. (2026). The Scale Fallacy: Why Semantic Gravity is Just a Bigger Hallucination. \emph{lab/sabine/colab/sabine\_the\_scale\_fallacy.tex}
\end{thebibliography}

\end{document}
